\section{相关工作}
\label{sec:related_work}

\subsection{果实检测方法}
果实检测方法可分为三类：基于2D图像的方法、基于3D点云的方法和多模态融合方法。

\textbf{基于2D图像的检测。} 基于RGB图像的果实检测已被广泛研究。传统方法依赖于HSV或Lab色彩空间中的颜色分割、形状分析（圆形度、椭圆度）和纹理特征提取。深度学习方法显著推动了领域发展，卷积神经网络（CNN）取得了较高的检测精度。DeepFruits系统使用深度神经网络在果园图像中实现了约85\%的苹果检测精度~\citep{Sa_2016}。后续工作通过改进网络架构和注意力机制提高了检测性能。CitrusYOLO通过改进YOLOv4并结合注意力模块，实现了96.15\%的柑橘检测平均精度~\citep{Chen_2022}。MangoDetNet引入了弱监督学习方法，在减少训练数据需求的同时保持检测性能~\citep{Denarda_2024}。

\textbf{基于3D点云的检测。} 三维传感器提供了能够进行果实尺寸估计和空间定位的几何信息。点云聚类算法，特别是DBSCAN及其变体，是分割识别3D空间中单个果实的主要方法~\citep{Ester_1996}。基于体积的估算方法已得到全面综述，比较了果实尺寸测量中从卡尺到机器视觉的方法~\citep{Neupane_2023}。深度信息与RGB数据的结合能够在不同光照条件下实现更稳健的检测。

\textbf{多模态融合。} 融合2D视觉信息与3D几何数据的方法解决了各单一模态的固有局限性。RGB-D相机和LiDAR传感器能够同时捕获颜色和深度信息。基于RGB-D的果实检测研究已证明通过互补的视觉和几何特征提高了稳健性~\citep{Tu_2020}。融合方法利用了2D检测的分类优势和3D点云处理的空间推理能力。

\subsection{点云聚类算法}
DBSCAN（基于密度聚类的噪声应用空间聚类）是一种基本的基于密度的聚类算法，它将密度连接的点分组为簇，同时自动识别噪声点~\citep{Ester_1996}。该算法需要两个参数：邻域半径$\epsilon$和核心点所需的最小点数MinPts。DBSCAN的优点包括无需预设簇数量、能够发现任意形状的簇以及对噪声的鲁棒性。

然而，标准DBSCAN使用固定的$\epsilon$参数，这对于LiDAR数据不是最优的，因为点密度随距传感器的距离显著变化。已经提出自适应DBSCAN变体来通过基于局部密度估计或距离相关计算动态调整邻域半径来解决这一限制。

\subsection{产量估算方法}
智能果实产量估算已在深度学习语义分割技术的背景下得到广泛综述~\citep{Maheswari_2021}。方法包括直接计数法（枚举检测到的果实并推断到树木或果园规模）、基于体积的估算（从3D重建计算果实体积并使用果实特定密度参数估算重量）以及冠层回归模型（将冠层特征与果实负载相关联）。

产量估算的准确性关键取决于遮挡问题的处理，因为在树冠内隐藏的果实在任何单一视角下都不可见。多角度扫描和统计校正方法已被提出用于解决这一挑战。